% Fragment: fig-sealed-residuals-converge.tex -- three residuals opened
% earlier in the chapter, one q*b account that prices all of them.
%
% Reader question: why do three separate gaps opened earlier in the
% chapter all get paid off by one number?
% Claim: the laundering residual, the malicious-worker gap in the ledger,
% and the arbitrary-telemetry channel are three faces of one debt, all
% priced by the same q*b capacity account -- not three independent
% problems needing three independent fixes.
% Evidence: opening paragraph of Sec. sealed-budget, website-v2/public/
% whitepaper/sealed-harbor.tex ("It is also where the chapter pays three
% debts it took on earlier: the laundering residual...the malicious
% worker...and the arbitrary-telemetry channel...All three are the same
% debt seen from three sides, and it has one price.").
% Must distinguish: three DIFFERENT sources (different sections, different
% failure shapes) from one SHARED sink -- the figure must not suggest
% three separate accounts that happen to look similar; there is exactly
% one q*b bucket, filled from three pipes.
% Grammar chosen: converging pipes into one bucket (S12's ledger
% convention: asset pools are buckets joined by directed pipes) -- reused
% here as three named residuals piping into the one leakage-budget bucket
% introduced in fig-sealed-leakage-ledger, this row's own paired panel.
% Grammar rejected: three separate small ledgers, one per residual --
% that is the "three independent problems" misreading the chapter is
% explicit about ruling out.
% Five-second test: count the buckets -- one, fed by three pipes, not
% three buckets side by side.
% QA: beauty_lint warns B7 in all three editions (sub-point
% near-alignment noise among the three stacked source nodes' centres);
% justified, not fixed.
\begin{figure}[H]
\centering
\pdfigurehue{pdcobalt}%
\begin{tikzpicture}[pd figure,x=1cm,y=1cm]
  % The explicit box leaves print-size air above the heavy first node. Without
  % it, fragment cropping trims to that node's inferred top rule.
  \path[use as bounding box] (-0.35,-1.35) rectangle (9.25,4.95);
  \def\dy{1.85}
  \node[pd title,anchor=south west] at (-0.10,4.45)
    {three residuals, one account};
  \node[pd breach state,minimum width=4.1cm,align=center] (R1) at (0,{2*\dy})
    {laundering residual\\\pdfigsub worker computes $t=f(s)$};
  \node[pd breach state,minimum width=4.1cm,align=center] (R2) at (0,{1*\dy})
    {malicious worker\\\pdfigsub declares its own $\varepsilon_i$};
  \node[pd breach state,minimum width=4.1cm,align=center] (R3) at (0,0)
    {arbitrary telemetry\\\pdfigsub unconstrained prose channel};
  \node[pd bucket,minimum width=3.0cm,minimum height=2.4cm,align=center]
    (ACCT) at (6.5,{1*\dy}) {\\[8pt]the $q\cdot b$\\account};
  \draw[pd focus pipe] (R1.east) -| (ACCT.north);
  \draw[pd focus pipe] (R2.east) -- (ACCT.west);
  \draw[pd focus pipe] (R3.east) -| (ACCT.south);
\end{tikzpicture}
\caption{Three named residuals converge on one leakage-capacity account. Each
lets the worker choose bits that the noninterference property does not forbid,
so all three are priced by the same $q\cdot b$ bound rather than by separate
mechanisms. [internal]}
\label{fig:sealed-residuals-converge}
\end{figure}
